\section{Glosario y lista de herramientas}
\subsection{Términos comunes}
Para facilitar la colaboración del equipo, Graham recomienda mantener un glosario: Plan-Act-Verify, Agent Chain, Observability Stream, entre otros; un lenguaje unificado reduce los malentendidos.

\begin{itemize}
    \item Plan-Act-Verify: divide con claridad los pasos de planeación, ejecución y verificación del agente.
    \item Agent Chain: combina varios agentes en un flujo de ejecución encadenado para tareas complejas.
    \item Observability Stream: el flujo de datos unificado de los registros, comandos y resultados de pruebas que produce el agente.
\end{itemize}

\subsection{Lista de herramientas}
OpenHands incorpora varios conjuntos de herramientas: Terminal Runner, Browser Controller, Test Executor, Diff Reporter. Al desplegar, se puede mantener un entorno y permisos independientes para cada una, y registrar su uso.

\begin{table}[H]
\centering
\begin{tabular}{p{0.3\textwidth} p{0.6\textwidth}}
\toprule
Herramienta & Uso \\
\midrule
Terminal Runner & Ejecuta comandos de shell y registra stdout/stderr. \\
\midrule
Browser Controller & Realiza interacción con páginas, extrae datos y los reescribe en el entorno. \\
\midrule
Test Executor & Ejecuta pytest, npm test, etcétera, y recopila los resultados. \\
\midrule
Diff Reporter & Reúne los patches y genera una descripción y un diff en markdown. \\
\bottomrule
\end{tabular}
\caption{Principales herramientas de OpenHands y su función.}
\end{table}

\begin{knowledgebox}{Estrategia de combinación de herramientas}
Procura usar herramientas con permisos distintos en pasos diferentes para reducir el riesgo de operaciones erróneas; por ejemplo, limita las operaciones de base de datos a un comando independiente y establece una lógica de reversión.
\end{knowledgebox}

\subsection{Resumen de la sección}
\begin{itemize}
    \item Un glosario facilita la colaboración entre equipos y evita malentendidos sobre la intención del agente.
    \item Una lista de herramientas ayuda a modularizar distintos tipos de operaciones, facilitando la asignación de permisos y la auditoría.
\end{itemize}
